% Pista alg-23j2-q5 · Àlgebra
% Procedència: PAU Catalunya 2023 · Convocatòria de juny · Sèrie 5 · Qüestió 5
\documentclass[11pt,a4paper]{article}
\usepackage{pista}
\pistainfo{Catalunya 2023}{Convocatòria de juny}{Sèrie 5}

\begin{document}

\section*{\textcolor{blauPista}{Qüestió 5}}

\noindent Considereu la família $\mathcal{S}$ de matrius de la forma $\begin{pmatrix} a & b \\ b & a \end{pmatrix}$, en què $a, b \in \mathbb{R}$.

\subsection*{\textit{a)} \normalfont Calculeu $\begin{pmatrix} 2 & 3 \\ 3 & 2 \end{pmatrix} \cdot \begin{pmatrix} 2 & 1 \\ 1 & 2 \end{pmatrix}^{-1}$. \hspace*{\fill}\textnormal{[1,25~punts]}}

\pista{En primer lloc, calcula la inversa de la segona matriu:
\[
\begin{pmatrix} 2 & 1 \\ 1 & 2 \end{pmatrix}
\]
Segueix els passos habituals: calcula el determinant, troba la transposada, troba l'adjunta, i divideix-ho tot pel determinant.

\medskip
Finalment, fes el producte: la primera matriu a l'esquerra i la inversa a la dreta.}

\subsection*{\textit{b)} \normalfont Trobeu totes les matrius de la família $\mathcal{S}$, és a dir, de la forma $A = \begin{pmatrix} a & b \\ b & a \end{pmatrix}$, que verifiquin la igualtat $A^{2} = I$, en què $I$ és la matriu identitat d'ordre 2. \hspace*{\fill}\textnormal{[1,25~punts]}}

\pista{En primer lloc, escriu la matriu $A$ amb les lletres $a$ i $b$, i calcula $A^{2} = A \cdot A$. Obtindràs una matriu que depèn d'uns paràmetres $a$ i $b$.

\medskip
A continuació, iguala la matriu obtinguda amb la matriu identitat $I$, i obtindràs quatre equacions. No és tan difícil com sembla, perquè algunes d'aquestes equacions surten repetides.

\medskip
Finalment, resol-les. Si una equació diu que un producte val zero, algun dels factors ha de valer zero: estudia cada possibilitat per separat. Al final, escriu les matrius, no només els valors dels paràmetres $a$ i $b$.}

\end{document}
